\chapter{Bloco 6 -- Problema da mochila binária}
\label{cap:bloco6}

Temos itens $j$ com \textbf{valor} $c_j$ e \textbf{peso} $a_j$, e uma mochila de capacidade $b$. Queremos o conjunto
de maior valor total que cabe. A mochila é um dos problemas clássicos mais importantes: aparece sozinha (orçamento
de projetos, carga de contêineres) e como \emph{subproblema} de modelos maiores.

\section{Quatro versões do problema}

\begin{center}
\begin{tabular}{ccccl}
\toprule
Caso & Fracionável? & Limite por item & Tipo & Solução\\
\midrule
1 & sim & ilimitado & PL & trivial: só o item de maior $c_j/a_j$, $x_{j^*}=b/a_{j^*}$\\
2 & não & ilimitado & PI & \textbf{não trivial}\\
3 & sim & $\le1$ & PL & gulosa por $c_j/a_j$, fracionando o último item\\
\textbf{4} & \textbf{não} & $\le\mathbf{1}$ & \textbf{PB} & \textbf{não trivial} $\leftarrow$ \emph{este bloco}\\
\bottomrule
\end{tabular}
\end{center}

\begin{exemplo}[caso 3 nos slides: a gulosa fracionária]
$\max 6x_1+8x_2+4x_3+x_4$ s.a.\ $x_1+2x_2+2x_3+x_4\le4$, $0\le x_j\le1$. Razões: $6$, $4$, $2$, $1$.
Pegamos $x_1=1$ (sobram 3), $x_2=1$ (sobra 1) e $x_3=1/2$ (acabou); $x_4=0$. Valor $6+8+2=16$.
\end{exemplo}

O caso 3 é exatamente a \textbf{relaxação linear} do caso 4. Então a gulosa fracionária dá, sem simplex nenhum, o
limite superior da mochila binária.

\begin{outros}[Dantzig (1957) e o ``item crítico'']
Esse resultado é clássico: Dantzig mostrou em 1957 que a relaxação contínua da mochila é resolvida ordenando os itens
por $c_j/a_j$ e enchendo a mochila até o primeiro item que não cabe inteiro, o \textbf{item crítico}, que entra
fracionado~\cite{dantzig1957,martellototh1990}. O valor resultante (o ``limite de Dantzig'') é o limite superior
usado nos primeiros algoritmos de \emph{branch-and-bound} para a mochila.
\end{outros}

\section{O problema do notebook: carteira de projetos}

Dez projetos, orçamento $B=74$ (R\$ mil):
\begin{center}
\small
\begin{tabular}{l*{10}{c}}
\toprule
projeto $j$ & 1&2&3&4&5&6&7&8&9&10\\
\midrule
valor $c_j$ & 29&47&14&33&37&13&41&36&38&16\\
custo $a_j$ & 34&25&19&32&18&12&16&36&16&16\\
$c_j/a_j$ & 0,85&1,88&0,74&1,03&2,06&1,08&\textbf{2,56}&1,00&2,38&1,00\\
\bottomrule
\end{tabular}
\end{center}

\begin{discussao}
\begin{enumerate}
  \item \emph{Variável? Por que binária e não inteira qualquer?} $x_j=1$ se o projeto $j$ é financiado. Um projeto
  não pode ser financiado ``duas vezes'' nem pela metade: é o caso 4.
  \item \emph{Objetivo e restrição?} $\max\sum_jc_jx_j$ s.a.\ $\sum_ja_jx_j\le74$. Uma única restrição!
  \item \emph{Qual entra primeiro pela razão? A gulosa resolve? E a relaxação, é o caso 3?} O projeto 7 (razão 2,56)
  entra primeiro. A gulosa \emph{não} garante o ótimo. A relaxação é o caso 3.
\end{enumerate}
\end{discussao}

\subsection*{Exercício 6A}
\begin{codigo}
def modelo_mochila(valor, peso, B, extras=False):
    prob = pulp.LpProblem("mochila", pulp.LpMaximize)
    x = {j: pulp.LpVariable(f"x_{j}", cat="Binary") for j in valor}
    prob += pulp.lpSum(valor[j]*x[j] for j in valor), "valor_total"       # TODO 1
    prob += pulp.lpSum(peso[j]*x[j] for j in valor) <= B, "capacidade"    # TODO 2
    if extras:
        ...                                                               # 6B, abaixo
    return prob, x
\end{codigo}
\py{valor} e \py{peso} são dicionários com chaves $1,\dots,10$ (\py{enumerate(..., start=1)}), então
\py{for j in valor} percorre os projetos.

\section{Comparando quatro respostas}

\paragraph{Gulosa (heurística).} Ordem por razão: 7, 9, 5, 2, 6, 4, 8, 10, 1, 3.
\begin{center}
\begin{tabular}{clccl}
\toprule
passo & projeto & custo acumulado & valor & decisão\\
\midrule
1 & 7 & 16 & 41 & entra\\
2 & 9 & 32 & 79 & entra\\
3 & 5 & 50 & 116 & entra\\
4 & 2 & 75 & -- & não cabe ($75>74$)\\
5 & 6 & 62 & 129 & entra\\
6+ & 4, 8, 10, 1, 3 & -- & -- & nenhum cabe nos 12 restantes\\
\bottomrule
\end{tabular}
\end{center}
Resultado: $\{5,6,7,9\}$, valor \textbf{129}.

\paragraph{Relaxação linear (Dantzig).} Mesma ordem, mas o projeto 2 entra \emph{fracionado} para usar os 24 que
sobram: $x_2=24/25=0{,}96$. Valor $116+0{,}96\cdot47=\mathbf{161{,}12}$. O projeto 2 é o item crítico.

\paragraph{Ótimo (PuLP/CBC) e programação dinâmica.} $\{2,7,9,10\}$, custo $25+16+16+16=73$, valor
$47+41+38+16=\mathbf{142}$. A programação dinâmica $O(nB)$ do notebook confirma (você verá PD mais adiante na
disciplina).

\begin{center}
\begin{tikzpicture}[x=0.075cm]
  \draw[->] (120,0) -- (168,0);
  \foreach \v in {120,130,140,150,160} \draw (\v,0.1) -- (\v,-0.1) node[below,font=\footnotesize] {\v};
  \fill[laranja] (129,0) circle (3pt) node[above=3pt,font=\small,align=center] {gulosa\\129};
  \fill[verde] (142,0) circle (3pt) node[above=3pt,font=\small,align=center] {ótimo\\142};
  \fill[azul] (161.12,0) circle (3pt) node[above=3pt,font=\small,align=center] {relaxação\\161,12};
  \draw[decorate,decoration={brace,mirror,amplitude=5pt}] (129,-0.6) -- (161.12,-0.6)
    node[midway,below=6pt,font=\small] {intervalo que o \emph{branch-and-bound} precisa fechar};
\end{tikzpicture}
\end{center}

\section{6B -- Mochila com regras de negócio}

Juntando os blocos 1 a 3:
\begin{center}
\begin{tabular}{lll}
\toprule
Regra & Lógica & Restrição\\
\midrule
(a) projetos 2 e 7 incompatíveis & exclusão mútua & $x_2+x_7\le1$\\
(b) se 9 então 5 & implicação & $x_9\le x_5$\\
(c) no máximo 4 projetos & cardinalidade & $\sum_jx_j\le4$\\
(d) 3 ou 6 (ou ambos) & ou inclusivo & $x_3+x_6\ge1$\\
\bottomrule
\end{tabular}
\end{center}
\begin{codigo}
    if extras:
        prob += x[2] + x[7] <= 1, "excl_2_7"          # (a)
        prob += x[9] <= x[5], "c9_implica_5"           # (b)
        prob += pulp.lpSum(x.values()) <= 4, "max4"   # (c)
        prob += x[3] + x[6] >= 1, "c3_ou_6"            # (d)
\end{codigo}
\paragraph{Resultado.} Ótimo \textbf{135} com $\{2,5,6,9\}$ (custo $25+18+12+16=71$). A regra (a) corta a solução de
6A (que tinha 2 e 7 juntos); (b) obriga o 5 a acompanhar o 9; (d) é satisfeita pelo 6.

\begin{atencao}[erros típicos]
\begin{itemize}
  \item Escrever (b) como $x_5\le x_9$: isso é ``se 5 então 9'', a implicação ao contrário.
  \item Escrever (d) como $x_3+x_6=1$: isso é o ``ou exclusivo'' e proíbe escolher os dois.
\end{itemize}
\end{atencao}

\begin{resumo}[Reflexão do notebook]
A gulosa fica em 129, a relaxação diz ``até 161'' e o ótimo é 142. Quem fecha a diferença é o
\emph{branch-and-bound}, que usa a relaxação como limite e ramifica nas variáveis fracionárias (aqui, primeiro em
$x_2$).
\end{resumo}
